\chapter{Conclusion and Future Work}

\section{Conclusion}

This thesis presented a technical study of an enterprise API connector development platform and confirmed that the implemented workflow is complete, practical, and maintainable for its target environment.

The platform unifies five phases---module management, contract parsing, EJB metadata extraction, mapping/config generation, and module generation---that were previously executed through disconnected tools. The result is improved consistency, lower onboarding effort, and traceable mapping decisions.

Its architecture combines a guided React frontend, an Express orchestration backend, parser/generator utility modules, and real-time feedback via Socket.IO. The hybrid synchronization model (watchers + events + polling fallback) provides resilient behavior under enterprise network constraints. The evaluation also identifies a clear enhancement path for security, portability, storage throughput, and testing depth.

\section{Key Contributions}

The thesis contributes:
\begin{enumerate}
    \item complete architectural documentation of frontend, backend, utility, storage, and generation layers;
    \item algorithm-level analysis of Swagger parsing (\texttt{\$ref} resolution, flattening, operation extraction);
    \item detailed documentation of JAR bytecode introspection (descriptor/signature parsing, EJB detection, deep path extraction, parameter normalization);
    \item evidence-based assessment across reliability, security, portability, and maintainability;
    \item a prioritised roadmap for progressive production hardening.
\end{enumerate}

\section{Future Work}

The platform is functionally mature, but further hardening is required for broader production use. Highest-priority work includes security controls, cross-platform generation support, storage evolution, and automated testing. Medium-priority work includes frontend refactoring and containerized deployment.

\noindent The roadmap is implemented as phased hardening:
\begin{itemize}
    \setlength\itemsep{2pt}
    \item \textbf{Security}: introduce JWT/RBAC, stricter sanitization, and tighter CORS/certificate policies for broader deployment boundaries.
    \item \textbf{Portability}: replace or wrap PowerShell generation with a cross-platform orchestrator while preserving Windows compatibility during transition.
    \item \textbf{Data and quality}: evolve storage for higher concurrency and add unit, integration, and end-to-end automated tests in CI.
    \item \textbf{Maintainability and scale}: decompose \texttt{YamlUploader}, add containerized deployment, and prepare for tenant isolation and queued generation workloads.
\end{itemize}

These enhancements extend the same core architecture and are feasible as incremental releases. The workflow model is also transferable to other regulated integration domains where API-to-enterprise mapping is required and source-code access may be limited.

\section{Summary}

The thesis shows that a workflow-driven connector platform is both technically feasible and operationally valuable in enterprise integration settings. The evaluated implementation delivers strong coverage across parsing, introspection, mapping, generation, synchronization, and deployment.

The roadmap defines incremental next steps that build directly on the current architecture, enabling hardening and scale-up without fundamental redesign.
